% ==========================================================
% Anexos: instrumentos de recolección de información y plantilla de matriz
% ==========================================================

\newcommand{\linea}[1]{\rule{#1}{0.4pt}}

% ----------------------------------------------------------
\Anexo{Lista de verificación de requisitos normativos}
\label{anx:lista_verificacion}

\textbf{Objetivo:} verificar, para el proceso analizado, el cumplimiento de los requisitos de la LOPDP, su Reglamento General y la normativa técnica de la SPDP que inciden en la gestión de riesgos (OE1). \textbf{Aplicación:} revisión documental y contraste con la práctica efectiva; para cada requisito se registra la evidencia observada, no solo la existencia de una política escrita. \textbf{Escala:} C = cumple; P = cumple parcialmente; N = no cumple; NA = no aplica.

\begin{center}
\footnotesize
\renewcommand{\arraystretch}{1.25}
\begin{tabularx}{\textwidth}{|c|L|p{2.6cm}|p{3.0cm}|c|}
  \hline
  \textbf{N.º} & \textbf{Requisito} & \textbf{Fundamento} & \textbf{Evidencia esperada} & \textbf{C/P/N/NA} \\
  \hline
  1 & El responsable mantiene un inventario o registro de las actividades de tratamiento del proceso. & LOPDP art.~40; guía SPDP (contexto) & RAT actualizado & \\
  \hline
  2 & Cada actividad de tratamiento tiene una base de legitimación identificada y documentada. & LOPDP, art.~7 & Matriz de bases de legitimación & \\
  \hline
  3 & El consentimiento, cuando es la base de legitimación, se solicita de forma diferenciada por finalidad. & LOPDP, art.~8 & Formularios y cláusulas & \\
  \hline
  4 & Se informa al titular sobre las finalidades, los destinatarios, los plazos y el ejercicio de sus derechos. & LOPDP, arts.~10, lit.~c, y 12 & Aviso de privacidad & \\
  \hline
  5 & Las categorías especiales de datos (salud, discapacidad, adolescentes) se identifican y reciben protección reforzada. & LOPDP art.~25 & Clasificación y controles específicos & \\
  \hline
  6 & Existen medidas de seguridad técnicas, organizativas y jurídicas, con verificación y evaluación continua de su eficacia. & LOPDP art.~37 & Políticas, registros de revisión & \\
  \hline
  7 & Se aplica una metodología de análisis de riesgos, amenazas y vulnerabilidades adaptada a la naturaleza del tratamiento. & LOPDP arts.~40 y 68 num.~4 & Metodología y matriz de riesgos & \\
  \hline
  8 & Se ha determinado si el tratamiento es de gran escala y, en su caso, se ha realizado la evaluación de impacto previa. & LOPDP art.~42; Res. SPDP-SPD-2026-0005-R & Dictamen de gran escala; EIPD & \\
  \hline
\end{tabularx}
\end{center}

\begin{center}
\footnotesize
\renewcommand{\arraystretch}{1.25}
\begin{tabularx}{\textwidth}{|c|L|p{2.6cm}|p{3.0cm}|c|}
  \hline
  \textbf{N.º} & \textbf{Requisito} & \textbf{Fundamento} & \textbf{Evidencia esperada} & \textbf{C/P/N/NA} \\
  \hline
  9 & Se ha designado un delegado de protección de datos cuando la normativa lo exige. & LOPDP, art.~48; Res. SPDP-SPD-2026-0005-R & Acto de designación & \\
  \hline
  10 & Existe un procedimiento para gestionar y notificar vulneraciones a la autoridad, al organismo de control y al CSIRT en el término de cinco días, y a los titulares en el de tres días, a través del SISPDP. & LOPDP, arts.~43 y 46; RLOPDP, arts.~24 a 28; Res. SPDP-SPD-2026-0040-R & Procedimiento y registro de incidentes & \\
  \hline
  11 & Los contratos con encargados (pasarela de pagos, nube, LMS, CRM) contienen cláusulas de protección de datos. & LOPDP, arts.~34 y 47, num.~10; RLOPDP, art.~41 & Contratos y anexos & \\
  \hline
  12 & Las transferencias a proveedores ubicados fuera del país cuentan con garantías adecuadas. & LOPDP, arts.~55 a 58; RLOPDP, arts.~71 a 78 & Evaluación de transferencias & \\
  \hline
  13 & Se han definido plazos de conservación y procedimientos de eliminación o anonimización. & LOPDP, art.~10, lit.~i; RLOPDP, art.~38 & Tabla de conservación; actas de eliminación & \\
  \hline
  14 & Existen canales y procedimientos para atender los derechos de los titulares en plazo. & LOPDP, arts.~12 a 24 & Procedimiento y registro de solicitudes & \\
  \hline
  15 & El personal que trata datos del proceso recibe capacitación periódica en protección de datos. & LOPDP, arts.~24 y 37 & Registros de capacitación & \\
  \hline
\end{tabularx}
\end{center}

\noindent\textbf{Nota:} los artículos citados se verificaron en el texto oficial de la LOPDP (Registro Oficial, 5.º Suplemento N.º 459) y del Reglamento General (Decreto Ejecutivo N.º 904), con la reforma del artículo 43 introducida por la LOFC (2026).

% ----------------------------------------------------------
\Anexo{Guía de entrevista semiestructurada}
\label{anx:entrevista}

\textbf{Objetivo:} caracterizar la práctica efectiva de tratamiento de datos personales en cada subproceso de admisión y matrícula (OE1). \textbf{Informantes:} responsables de las unidades que intervienen en el proceso de admisión y matrícula de la UTI (Admisiones, Secretaría Académica, Bienestar Estudiantil, Dirección Financiera y Dirección de TI, o sus equivalentes). \textbf{Duración estimada:} 45 a 60 minutos. \textbf{Registro:} notas del entrevistador y, con autorización, grabación de audio que se elimina tras la transcripción. Los entrevistados se identifican con códigos (E01, E02, \ldots).

\subsection*{Protocolo de inicio}
\begin{enumerate}
  \item Presentar el propósito académico del estudio y la autorización institucional con la que se realiza.
  \item Confirmar la firma del consentimiento informado (Anexo~\ref{anx:consentimiento}) y la autorización de grabación.
  \item Recordar que no se solicitarán datos personales de estudiantes ni de aspirantes.
\end{enumerate}

\subsection*{Bloque 1. Finalidades y bases de legitimación}
\begin{enumerate}
  \item ¿Qué datos personales se recolectan en su unidad durante el proceso de admisión y matrícula y con qué finalidad?
  \item ¿Cómo se informa al aspirante sobre el uso de sus datos? ¿Se solicita su consentimiento? ¿Para qué finalidades?
  \item ¿Se utilizan los datos para fines distintos de la admisión (difusión, marketing, investigación)?
\end{enumerate}

\subsection*{Bloque 2. Categorías especiales y grupos vulnerables}
\begin{enumerate}
  \item ¿Se tratan datos de salud, discapacidad, condición socioeconómica o de menores de edad? ¿En qué subproceso?
  \item ¿Qué tratamiento diferenciado reciben esos datos (acceso, almacenamiento, conservación)?
\end{enumerate}

\subsection*{Bloque 3. Sistemas, accesos y terceros}
\begin{enumerate}
  \item ¿En qué sistemas, repositorios o servicios en la nube se almacenan los datos?
  \item ¿Cómo se asignan, revisan y retiran los permisos de acceso? ¿Qué ocurre con el personal temporal?
  \item ¿Qué proveedores externos acceden a los datos? ¿Existen cláusulas contractuales de protección de datos?
\end{enumerate}

\subsection*{Bloque 4. Conservación y eliminación}
\begin{enumerate}
  \item ¿Cuánto tiempo se conservan los expedientes de aspirantes admitidos y no admitidos?
  \item ¿Existe un procedimiento de eliminación o anonimización? ¿Quién lo ejecuta y cómo se registra?
\end{enumerate}

\subsection*{Bloque 5. Incidentes y riesgos}
\begin{enumerate}
  \item ¿Conoce incidentes de seguridad o de exposición de datos ocurridos en el proceso (sin identificar personas)?
  \item ¿Qué haría su unidad ante la pérdida o exposición de datos? ¿A quién reportaría y en qué plazo?
  \item ¿Cuáles considera que son los tres principales riesgos para los titulares en su subproceso?
\end{enumerate}

\subsection*{Cierre}
Agradecer la participación, preguntar si desea agregar algo y explicar el uso que se dará a la información.

% ----------------------------------------------------------
\Anexo{Cuestionario de madurez en protección de datos personales}
\label{anx:madurez}

\textbf{Objetivo:} establecer la línea base de madurez de la institución en la gestión de datos personales del proceso analizado (fase 0, OE1). \textbf{Instrucciones:} para cada enunciado, seleccione el nivel que mejor describe la situación \textit{real} de la institución, no la deseada ni la escrita en una política. El nivel de cada dominio es la mediana de sus enunciados; el nivel global es la mediana de los dominios.

\begin{center}
\small
\renewcommand{\arraystretch}{1.25}
\begin{tabularx}{\textwidth}{|c|p{2.4cm}|L|}
  \hline
  \textbf{Nivel} & \textbf{Denominación} & \textbf{Descripción} \\
  \hline
  1 & Inicial & No existen prácticas definidas; las actividades dependen de iniciativas individuales. \\
  \hline
  2 & Repetible & Existen prácticas informales que se repiten, sin documentación ni responsables formales. \\
  \hline
  3 & Definido & Las prácticas están documentadas, aprobadas y asignadas a responsables. \\
  \hline
  4 & Gestionado & Las prácticas se miden con indicadores y se revisan periódicamente. \\
  \hline
  5 & Optimizado & Las prácticas se mejoran de forma continua a partir de los resultados y de los incidentes. \\
  \hline
\end{tabularx}
\end{center}

\begin{center}
\footnotesize
\renewcommand{\arraystretch}{1.2}
\begin{tabularx}{\textwidth}{|p{2.7cm}|c|L|c|c|c|c|c|}
  \hline
  \textbf{Dominio} & \textbf{N.º} & \textbf{Enunciado} & \textbf{1} & \textbf{2} & \textbf{3} & \textbf{4} & \textbf{5} \\
  \hline
  \multirow{3}{2.7cm}{Gobernanza} & 1 & Existe una política de protección de datos aprobada por la autoridad institucional. & & & & & \\
  \cline{2-8}
   & 2 & Hay un responsable o delegado de protección de datos con funciones definidas. & & & & & \\
  \cline{2-8}
   & 3 & La protección de datos se considera en la planificación de nuevos procesos y sistemas. & & & & & \\
  \hline
  \multirow{3}{2.7cm}{Inventario y licitud} & 4 & Se mantiene un RAT actualizado del proceso. & & & & & \\
  \cline{2-8}
   & 5 & Cada tratamiento tiene una base de legitimación documentada. & & & & & \\
  \cline{2-8}
   & 6 & Los formularios informan al titular y separan las finalidades. & & & & & \\
  \hline
  \multirow{3}{2.7cm}{Gestión de riesgos} & 7 & Se identifican y evalúan los riesgos para los titulares. & & & & & \\
  \cline{2-8}
   & 8 & Se determina cuándo se requiere una EIPD. & & & & & \\
  \cline{2-8}
   & 9 & Los riesgos tienen responsables, planes de tratamiento y seguimiento. & & & & & \\
  \hline
  \multirow{3}{2.7cm}{Seguridad} & 10 & Los accesos se asignan por función y se revisan periódicamente. & & & & & \\
  \cline{2-8}
   & 11 & Los datos sensibles se clasifican, cifran y respaldan. & & & & & \\
  \cline{2-8}
   & 12 & Se registran y monitorean los accesos a datos sensibles. & & & & & \\
  \hline
  \multirow{2}{2.7cm}{Terceros} & 13 & Los contratos con encargados incluyen cláusulas de protección de datos. & & & & & \\
  \cline{2-8}
   & 14 & Se evalúan las transferencias a proveedores en la nube. & & & & & \\
  \hline
  \multirow{2}{2.7cm}{Ciclo de vida} & 15 & Existen plazos de conservación definidos. & & & & & \\
  \cline{2-8}
   & 16 & Los datos se eliminan o anonimizan al vencer el plazo. & & & & & \\
  \hline
  \multirow{2}{2.7cm}{Derechos e incidentes} & 17 & Las solicitudes de los titulares se atienden en plazo y se registran. & & & & & \\
  \cline{2-8}
   & 18 & Existe un procedimiento probado de notificación de vulneraciones. & & & & & \\
  \hline
  \multirow{2}{2.7cm}{Cultura} & 19 & El personal recibe capacitación periódica en protección de datos. & & & & & \\
  \cline{2-8}
   & 20 & Los incidentes se analizan para mejorar los controles. & & & & & \\
  \hline
\end{tabularx}
\end{center}

% ----------------------------------------------------------
\Anexo{Cuestionario Delphi de estimación de riesgos}
\label{anx:delphi}

\textbf{Objetivo:} estimar, mediante consenso de expertos, la probabilidad anual y el impacto de cada escenario de riesgo del caso de estudio (OE3; su validez se evalúa en el OE4). \textbf{Participantes:} panel de al menos cinco expertos, identificados con códigos. \textbf{Procedimiento:} hasta tres rondas anónimas; desde la segunda ronda, cada experto recibe la mediana, el rango intercuartílico (RIC) y los argumentos anonimizados del panel y puede revisar su estimación. Se considera alcanzado el consenso cuando el RIC es igual o menor que 1 y al menos el 70\,\% de los expertos se sitúa a no más de un nivel de la mediana; en cada ronda se calcula además la \(W\) de Kendall. Las respuestas se registran en la hoja \textit{Delphi} de la plantilla de matriz (Anexo~\ref{anx:plantilla_matriz}).

\textbf{Escalas:} probabilidad anual e impacto sobre el titular según las Tablas~\ref{tab:escala_probabilidad} y~\ref{tab:escala_impacto} del Capítulo~\ref{chap:metodologia} (1 = muy baja/muy bajo; 5 = muy alta/muy alto).

\vspace{6pt}
\noindent Código del experto: \linea{2cm} \quad Ronda: \linea{1.5cm} \quad Fecha: \linea{3cm}

\begin{center}
\small
\renewcommand{\arraystretch}{1.6}
\begin{tabularx}{\textwidth}{|>{\raggedright\arraybackslash}p{3.4cm}|L|}
  \hline
  \multicolumn{2}{|l|}{\textbf{Escenario} \linea{1.5cm} \quad \textbf{Dimensión:} C \; I \; D \; Licitud} \\
  \hline
  Descripción del escenario & (comunidad de amenaza, vector, vulnerabilidad, datos y titulares afectados) \\
  \hline
  Resultado de la ronda anterior & Mediana: \linea{1cm} \quad RIC: \linea{1cm} \quad Argumentos del panel: \\
  \hline
  Probabilidad anual (1--5) & \linea{1cm} \\
  \hline
  Justificación de la probabilidad & \\[18pt]
  \hline
  Impacto sobre el titular (1--5) & \linea{1cm} \\
  \hline
  Justificación del impacto & \\[18pt]
  \hline
  Grado de confianza en la estimación & Bajo \quad Medio \quad Alto \\
  \hline
\end{tabularx}
\end{center}

% ----------------------------------------------------------
\Anexo{Afirmaciones de calibración de expertos}
\label{anx:calibracion}

\textbf{Objetivo:} medir la calibración de cada experto antes de la primera ronda Delphi mediante el puntaje de Brier (Ecuación~\ref{eq:brier}). \textbf{Instrucciones:} indique la probabilidad, entre 0 y 1, de que cada afirmación sea verdadera. No consulte fuentes; responda con su conocimiento actual. Las respuestas se registran en la hoja \textit{Brier} de la plantilla (Anexo~\ref{anx:plantilla_matriz}); los resultados reales se revelan al finalizar.

\begin{center}
\small
\renewcommand{\arraystretch}{1.3}
\begin{tabularx}{\textwidth}{|c|L|p{2.4cm}|}
  \hline
  \textbf{N.º} & \textbf{Afirmación} & \textbf{Probabilidad de que sea verdadera (0--1)} \\
  \hline
  C01 & La LOPDP se publicó en el Registro Oficial en 2021. & \\
  \hline
  C02 & El artículo 43 de la LOPDP fija un término de cinco días para notificar vulneraciones a la autoridad. & \\
  \hline
  C03 & ISO/IEC 27005:2022 es una norma certificable. & \\
  \hline
  C04 & El Reglamento General a la LOPDP se expidió mediante Decreto Ejecutivo N.º 904. & \\
  \hline
  C05 & Según el estudio de Sophos de 2024, más del 60\,\% de las IES encuestadas sufrió un ataque de \textit{ransomware} en el último año. & \\
  \hline
  C06 & ISO/IEC 27002:2022 agrupa sus controles en cuatro temas. & \\
  \hline
  C07 & ISO/IEC 27002:2022 contiene más de 100 controles. & \\
  \hline
  C08 & El Reglamento General de Protección de Datos europeo es aplicable desde mayo de 2018. & \\
  \hline
  C09 & El marco LINDDUN incluye la categoría de amenaza \textit{spoofing}. & \\
  \hline
  C10 & La Resolución SPDP-SPD-2026-0005-R regula el tratamiento de datos personales a gran escala. & \\
  \hline
\end{tabularx}
\end{center}

% ----------------------------------------------------------
\Anexo{Autoevaluación del coeficiente de competencia experta}
\label{anx:competencia}

\textbf{Objetivo:} determinar el coeficiente de competencia $K = \frac{1}{2}(K_c + K_a)$ de cada candidato al panel (Ecuación~\ref{eq:competencia}), conforme al procedimiento descrito por Cabero Almenara y Barroso Osuna \cite{cabero2013competencia}. \textbf{Instrucciones:} responda las dos partes con base en su experiencia real; sus respuestas se identifican solo con su código.

\vspace{6pt}
\noindent Código del candidato: \linea{2cm} \quad Perfil profesional: \linea{4.5cm} \quad Años de experiencia: \linea{1.5cm}

\vspace{6pt}
\noindent\textbf{Parte 1. Coeficiente de conocimiento ($K_c$).} Marque con una X su nivel de conocimiento sobre gestión de riesgos en el tratamiento de datos personales en instituciones de educación superior, en una escala de 0 (ningún conocimiento) a 10 (pleno conocimiento). $K_c$ es el valor marcado multiplicado por 0,1.

\begin{center}
\small
\renewcommand{\arraystretch}{1.4}
\begin{tabular}{|c|c|c|c|c|c|c|c|c|c|c|}
  \hline
  0 & 1 & 2 & 3 & 4 & 5 & 6 & 7 & 8 & 9 & 10 \\
  \hline
  & & & & & & & & & & \\
  \hline
\end{tabular}
\end{center}

\noindent\textbf{Parte 2. Coeficiente de argumentación ($K_a$).} Indique el grado de influencia de cada fuente en su criterio. $K_a$ es la suma de los pesos de las casillas marcadas.

\begin{center}
\small
\renewcommand{\arraystretch}{1.3}
\begin{tabularx}{\textwidth}{|L|c|c|c|}
  \hline
  \textbf{Fuente de argumentación} & \textbf{Alto} & \textbf{Medio} & \textbf{Bajo} \\
  \hline
  Análisis teóricos realizados por usted & 0,3 & 0,2 & 0,1 \\
  \hline
  Experiencia práctica obtenida & 0,5 & 0,4 & 0,2 \\
  \hline
  Trabajos de autores nacionales & 0,05 & 0,05 & 0,05 \\
  \hline
  Trabajos de autores extranjeros & 0,05 & 0,05 & 0,05 \\
  \hline
  Conocimiento del estado del tema en el extranjero & 0,05 & 0,05 & 0,05 \\
  \hline
  Su intuición & 0,05 & 0,05 & 0,05 \\
  \hline
\end{tabularx}
\end{center}

\noindent\textbf{Interpretación:} competencia alta si $K > 0{,}8$; media si $0{,}5 \le K \le 0{,}8$; baja si $K < 0{,}5$ \cite{cabero2013competencia}.

% ----------------------------------------------------------
\Anexo{Instrumento de validación por juicio de expertos}
\label{anx:validacion}

\textbf{Objetivo:} valorar la metodología propuesta y sus productos (OE4). \textbf{Instrucciones:} tras revisar el documento metodológico, la plantilla de matriz y el expediente anonimizado del caso de la UTI, califique cada criterio en una escala de 1 (totalmente en desacuerdo) a 5 (totalmente de acuerdo) y registre sus observaciones. Para cada enunciado se calcula la mediana y la V de Aiken (Ecuación~\ref{eq:aiken}); un enunciado se considera validado cuando \(V \ge 0{,}70\) y la mediana es igual o mayor que 4, y un criterio, cuando lo están todos sus enunciados.

\vspace{6pt}
\noindent Código del experto: \linea{2cm} \quad Perfil profesional: \linea{4.5cm} \quad Años de experiencia: \linea{1.5cm}

\begin{center}
\footnotesize
\renewcommand{\arraystretch}{1.25}
\begin{tabularx}{\textwidth}{|c|p{2.3cm}|L|c|c|c|c|c|}
  \hline
  \textbf{N.º} & \textbf{Criterio} & \textbf{Enunciado} & \textbf{1} & \textbf{2} & \textbf{3} & \textbf{4} & \textbf{5} \\
  \hline
  V01 & Pertinencia & La metodología responde a las obligaciones de análisis de riesgos y evaluación de impacto de la LOPDP. & & & & & \\
  \hline
  V02 & Pertinencia & Es adecuada para el contexto de una IES ecuatoriana. & & & & & \\
  \hline
  V03 & Coherencia & Sus fases siguen las etapas de la guía de la SPDP. & & & & & \\
  \hline
  V04 & Coherencia & Las estimaciones exigen justificación y se refieren a un lapso anual. & & & & & \\
  \hline
  V05 & Claridad & Las actividades y productos de cada fase se entienden sin ambigüedad. & & & & & \\
  \hline
  V06 & Claridad & Las escalas de probabilidad e impacto son comprensibles y discriminan entre niveles. & & & & & \\
  \hline
  V07 & Suficiencia & La valoración del impacto refleja el daño a los titulares y no solo a la institución. & & & & & \\
  \hline
  V08 & Suficiencia & La tabla de decisión y los criterios de aceptación tratan adecuadamente los riesgos de baja probabilidad y alto impacto. & & & & & \\
  \hline
  V09 & Pertinencia & La metodología puede aplicarse con los recursos habituales de una IES. & & & & & \\
  \hline
  V10 & Claridad & La plantilla de matriz es fácil de usar y reutilizable en otros procesos. & & & & & \\
  \hline
  V11 & Pertinencia & Los resultados permiten priorizar controles y documentar la responsabilidad proactiva. & & & & & \\
  \hline
  V12 & Suficiencia & La fase de monitoreo permite mantener la gestión de riesgos actualizada. & & & & & \\
  \hline
\end{tabularx}
\end{center}

\noindent\textbf{Preguntas abiertas}
\begin{enumerate}
  \item ¿Qué elementos de la metodología considera más valiosos?
  \item ¿Qué elementos modificaría, eliminaría o agregaría?
  \item ¿Qué obstáculos prevé para su adopción en una IES?
\end{enumerate}

% ----------------------------------------------------------
\Anexo{Protocolo del simulacro de vulneración (escenario de validación V2)}
\label{anx:simulacro}

\textbf{Objetivo:} comprobar que la fase de preparación y el procedimiento de respuesta de la metodología permiten clasificar, notificar y comunicar una vulneración dentro de los plazos legales (indicadores IV10 e IV11, OE4). \textbf{Modalidad:} ejercicio de mesa, sin intervenir sistemas reales ni usar datos personales reales. \textbf{Participantes:} responsables de Bienestar Estudiantil, Dirección de TI, área jurídica y, si existe, el delegado de protección de datos, identificados con códigos. \textbf{Duración:} dos a tres horas. \textbf{Escenario:} R-09, exfiltración de expedientes de becas en un ataque de doble extorsión; todos los volúmenes y datos del enunciado son ficticios.

\begin{center}
\small
\renewcommand{\arraystretch}{1.3}
\begin{tabularx}{\textwidth}{|c|>{\raggedright\arraybackslash}p{3.4cm}|L|c|}
  \hline
  \textbf{Hito} & \textbf{Inyección del facilitador} & \textbf{Decisión o producto esperado} & \textbf{Hora registrada} \\
  \hline
  T0 & Alerta de actividad anómala en el repositorio de expedientes de becas & Registro del incidente con hora de detección (inicio del cómputo) & \\
  \hline
  T1 & Confirmación de que el repositorio contiene datos socioeconómicos y certificados de discapacidad & Evaluación inicial; verificación de las causales del art.~24 del RLOPDP & \\
  \hline
  T2 & Mensaje de extorsión con muestra de archivos & Clasificación por dimensión (confidencialidad) y carácter (temporal, permanente o irreversible); valoración del impacto con la escala de la metodología & \\
  \hline
  T3 & Solicitud de la dirección de un informe para decidir & Borrador de notificación con los ocho elementos del art.~26 del RLOPDP, para el SISPDP, el organismo de control y el CSIRT & \\
  \hline
  T4 & Pregunta sobre la comunicación a los titulares & Decisión motivada sobre la comunicación (art.~46 de la LOPDP) y borrador en lenguaje claro & \\
  \hline
  T5 & Cierre del ejercicio & Lecciones aprendidas y actualización propuesta de la matriz & \\
  \hline
\end{tabularx}
\end{center}

\noindent\textbf{Lista de comprobación de la notificación (art.~26 del RLOPDP):} \(\square\) naturaleza y tipo de vulneración; \(\square\) titulares afectados; \(\square\) detalle inicial de los sistemas vulnerados; \(\square\) causa presunta; \(\square\) volumen y tipos de datos comprometidos; \(\square\) medidas adoptadas y previstas; \(\square\) evaluación del riesgo para los titulares; \(\square\) otros aspectos requeridos por la autoridad.

\noindent\textbf{Criterios de éxito:} IV10, tiempo entre T0 y el borrador de T3 menor o igual que 72 horas simuladas (meta interna) y, en ningún caso, superior al término de cinco días; IV11, los ocho elementos presentes en el borrador.

% ----------------------------------------------------------
\Anexo{Consentimiento informado}
\label{anx:consentimiento}

\begin{center}
\textbf{CONSENTIMIENTO INFORMADO PARA PARTICIPAR EN EL ESTUDIO}\\
\textit{Metodología para el diagnóstico y gestión de riesgos en el tratamiento de datos personales}
\end{center}

\textbf{Investigador:} Ramón Alejandro Manzaba Castro, estudiante de Ingeniería en Tecnologías de la Información, Universidad Tecnológica Indoamérica. \textbf{Tutora:} María Cristina Jiménez Hernández, MSc.

\textbf{Propósito.} Usted ha sido invitado a participar como experto en un trabajo de integración curricular cuyo propósito es diseñar y validar una metodología de gestión de riesgos para el tratamiento de datos personales en instituciones de educación superior. El caso de estudio es el proceso de admisión y matrícula de la Universidad Tecnológica Indoamérica, sede Quito; no se recogerán datos personales de estudiantes ni de aspirantes.

\textbf{Participación.} Su participación consiste en responder, según corresponda, una entrevista, un cuestionario de calibración, hasta tres rondas de un cuestionario Delphi y un instrumento de validación. El tiempo estimado total es de tres horas, distribuidas en varias semanas.

\textbf{Datos que se recogerán.} Sus respuestas, su perfil profesional y sus años de experiencia. Su nombre no aparecerá en el documento; se le asignará un código (E01, E02, \ldots). Si la entrevista se graba, el audio se eliminará después de su transcripción.

\textbf{Finalidad y conservación.} La información se utilizará exclusivamente con fines académicos para este trabajo, se almacenará en medios con acceso restringido y se eliminará al concluir el proceso de titulación, salvo las respuestas anonimizadas que forman parte del documento final.

\textbf{Voluntariedad y derechos.} Su participación es voluntaria y no recibirá compensación. Puede retirarse en cualquier momento sin necesidad de justificarlo y solicitar el acceso, la rectificación o la eliminación de sus respuestas escribiendo al investigador, conforme a la Ley Orgánica de Protección de Datos Personales.

\textbf{Riesgos y beneficios.} La participación no supone riesgos previsibles. El beneficio esperado es contribuir a mejorar la protección de los datos personales en las instituciones de educación superior.

\textbf{Contacto.} ramonalejandromanzaba@gmail.com

\vspace{10pt}
He leído la información anterior, he podido hacer preguntas y acepto participar voluntariamente.

\vspace{8pt}
\begin{tabular}{@{}ll@{}}
Código asignado: & \linea{4cm} \\[12pt]
Autorizo la grabación de audio de la entrevista: & Sí \quad No \\[12pt]
Firma del participante: & \linea{5cm} \\[12pt]
Firma del investigador: & \linea{5cm} \\[12pt]
Fecha: & \linea{4cm} \\
\end{tabular}

% ----------------------------------------------------------
\Anexo{Plantilla de matriz de riesgos}
\label{anx:plantilla_matriz}

La plantilla reutilizable de la matriz de riesgos (OE3) se entrega como archivo de hoja de cálculo \texttt{matriz/Plantilla\_Matriz\_Riesgos\_LOPDP.xlsx}, compatible con Microsoft Excel y LibreOffice Calc. Su estructura implementa el diseño de la Sección~\ref{sec:diseno_solucion}. Las celdas de ingreso se distinguen de las celdas con fórmula, y los campos categóricos usan listas desplegables. Los niveles de riesgo se colorean con formato condicional. La Tabla~\ref{tab:hojas_plantilla} describe sus hojas.

\begin{table}[htbp]
  \centering
  \caption{Hojas de la plantilla de matriz de riesgos.}
  \label{tab:hojas_plantilla}
  \footnotesize
  \renewcommand{\arraystretch}{1.25}
  \begin{tabularx}{\textwidth}{|p{2.6cm}|L|}
    \hline
    \textbf{Hoja} & \textbf{Contenido} \\
    \hline
    Instrucciones & Propósito, leyenda de colores, pasos por fase y reglas de la guía de la SPDP. \\
    \hline
    Matriz & Una fila por escenario y dimensión, con 34 campos agrupados en seis bloques. Calcula el impacto ajustado, el nivel inherente y el residual, la necesidad de EIPD y el criterio de aceptación. \\
    \hline
    Decisión & Tabla de decisión del nivel de riesgo y criterios de aceptación. \\
    \hline
    Escalas & Escalas de probabilidad anual, de impacto sobre el titular y sobre el responsable, y criterios agravantes. \\
    \hline
    Amenazas & Catálogo inicial de doce amenazas del proceso de admisión y matrícula. \\
    \hline
    Controles & Catálogo de controles candidatos con referencia a ISO/IEC 27002:2022. \\
    \hline
    Delphi & Registro de rondas con mediana, cuartiles, RIC, consenso y media ponderada por calibración. \\
    \hline
    Brier & Afirmaciones de calibración, puntaje de Brier y peso de cada experto. \\
    \hline
    Mapa de calor & Número de escenarios por probabilidad e impacto (inherente y residual) y resumen por nivel. \\
    \hline
    Listas & Valores de las listas desplegables. \\
    \hline
  \end{tabularx}
  \raggedright
  \elaboradopor{Ramón Alejandro Manzaba Castro}
  \fuente{Elaboración propia}
\end{table}

% ----------------------------------------------------------
\Anexo{Matriz de consistencia}
\label{anx:consistencia}

La Tabla~\ref{tab:matriz_consistencia} verifica la correspondencia entre el problema, los objetivos, la metodología, los resultados y las conclusiones del trabajo.

\begin{table}[htbp]
  \centering
  \caption{Matriz de consistencia del trabajo.}
  \label{tab:matriz_consistencia}
  \scriptsize
  \renewcommand{\arraystretch}{1.2}
  \begin{tabularx}{\textwidth}{|>{\raggedright\arraybackslash}p{2.6cm}|L|L|L|>{\raggedright\arraybackslash}p{2.0cm}|}
    \hline
    \textbf{Objetivo} & \textbf{Pregunta directriz} & \textbf{Metodología (técnicas e instrumentos)} & \textbf{Resultado y ubicación} & \textbf{Conclusión} \\
    \hline
    \multicolumn{5}{|>{\raggedright\arraybackslash}p{\dimexpr\textwidth-2\tabcolsep\relax}|}{\textbf{Problema:} ¿Cómo diagnosticar y gestionar de forma sistemática los riesgos del tratamiento de datos personales de estudiantes y docentes en una IES, de modo que se cumplan las exigencias de la LOPDP y de la normativa técnica de la SPDP? \textbf{Objetivo general:} diseñar la metodología (Capítulo~\ref{chap:introduccion}).} \\
    \hline
    OE1. Diagnóstico & ¿Qué se necesita para diagnosticar y qué línea base arroja el diagnóstico? & Revisión documental, entrevistas, cuestionario de madurez, auditoría de procesos (Anexos~\ref{anx:lista_verificacion}--\ref{anx:madurez}) & Requisitos D01--D10 (Tabla~\ref{tab:requisitos_diagnostico}); madurez y verificación normativa (Tablas~\ref{tab:madurez_resultados} y~\ref{tab:verificacion_normativa}, pendientes) & Capítulo~\ref{chap:conclusiones}, OE1 \\
    \hline
    OE2. Universidades & ¿En qué situación están las universidades y qué referentes orientan la metodología? & Análisis documental comparado de fuentes oficiales y estudios indexados & Cifras del sector (Tabla~\ref{tab:cifras_sector}); comparación de países (Tabla~\ref{tab:comparativa_paises}); decisiones de diseño (Tabla~\ref{tab:insumos_oe2}) & Capítulo~\ref{chap:conclusiones}, OE2 \\
    \hline
    OE3. Implementación & ¿Qué escenarios, niveles de riesgo y medidas resultan de implementar la metodología? & Fases P a 4; panel Delphi; plantilla de matriz (Anexos~\ref{anx:delphi} y~\ref{anx:plantilla_matriz}) & Escenarios (Tabla~\ref{tab:escenarios}); riesgo inherente y residual (Tablas~\ref{tab:riesgo_inherente} y~\ref{tab:riesgo_residual}, pendientes); trazabilidad (Tabla~\ref{tab:trazabilidad}) & Capítulo~\ref{chap:conclusiones}, OE3 \\
    \hline
    OE4. Evaluación & ¿Qué validez y aplicabilidad tiene la metodología? & Pruebas técnicas; calibración (Brier) y competencia ($K$); Delphi; juicio de expertos (V de Aiken); simulacro (Anexos~\ref{anx:calibracion}--\ref{anx:simulacro}) & Indicadores IV01--IV12 (Tabla~\ref{tab:indicadores_validacion}; IV01--IV03 cumplidos, IV04--IV12 pendientes) & Capítulo~\ref{chap:conclusiones}, OE4 \\
    \hline
  \end{tabularx}
  \raggedright
  \elaboradopor{Ramón Alejandro Manzaba Castro}
  \fuente{Elaboración propia}
\end{table}

% ----------------------------------------------------------
\Anexo{Autorización institucional para el estudio}
\label{anx:autorizacion}

\noindent\pend{oficio de autorización de la Universidad Tecnológica Indoamérica para realizar el estudio en el proceso de admisión y matrícula de la sede Quito. Debe indicar la instancia que autoriza, la fecha, el alcance (unidades, documentos e informantes a los que se accede), las condiciones de confidencialidad y anonimización, y el enlace institucional designado. Se incorpora escaneado una vez emitido, sin datos personales de terceros.}

\vspace{8pt}
\noindent\textbf{Condiciones del estudio que la autorización debe cubrir:} (a) acceso a documentos institucionales del proceso, en forma anonimizada o agregada; (b) entrevistas a los responsables de Admisiones, Secretaría Académica, Bienestar Estudiantil, Dirección Financiera, Dirección de TI y área jurídica; (c) datos agregados sobre el volumen de aspirantes y la proporción de menores de edad; (d) realización de un simulacro de mesa sin intervención en sistemas reales, y (e) publicación de los resultados en el trabajo de integración curricular en términos de procesos y controles, sin identificar a personas.
